% Assignment source: ne630/homework/markdown/HW11.md.
% Legacy foundation: ne630_problems/lesson_11/statement.tex and
% OpenMC_Problem_HW11.ipynb.

\noindent{\bf Problem 1 [FNRP 3.9]}. Suppose that the Maxwell-Boltzmann
distribution, Eq.~(2.34), represents the neutron density in Eqs.~(3.43) and
(3.44).
\begin{enumerate}[label=(\alph*)]
\item Find the value of $\bar{v}$.
\item If we define $\bar{E}=\frac{1}{2}m\bar{v}^2$, show that
      $\bar{E}\approx1.273kT$.
\item Why is your result different from the average energy of
      $\frac{3}{2}kT$ given by Eq.~(2.33)?
\end{enumerate}

\begin{adaptedsolution}
The Maxwell-Boltzmann energy distribution is
\[
 M(E)=\frac{2\pi}{(\pi kT)^{3/2}}E^{1/2}e^{-E/kT}.
\]
\begin{enumerate}[label=(\alph*)]
\item Substitution into Eq.~(3.44) gives
\begin{align*}
 \bar{v}
 &=\frac{\int_0^\infty v(E)M(E)\,dE}{\int_0^\infty M(E)\,dE}\\
 &=\frac{2\pi}{(\pi kT)^{3/2}}\sqrt{\frac{2}{m}}
   \int_0^\infty E e^{-E/kT}\,dE\\
 &=\boxed{\frac{2\sqrt{2kT}}{\sqrt{\pi m}}}.
\end{align*}
\item Therefore,
\[
 \bar{E}=\frac{1}{2}m\bar{v}^2
 =\frac{1}{2}m\left(\frac{2\sqrt{2kT}}{\sqrt{\pi m}}\right)^2
 =\boxed{\frac{4kT}{\pi}\approx1.273kT}.
\]
\item The results differ because $\bar{E}$ here is computed from the square of
the mean speed, while the true expected energy depends on the mean of the
speed squared.  In general, $\left(\bar{v}\right)^2\ne\overline{v^2}$.
\end{enumerate}
\end{adaptedsolution}

\newpage

\noindent{\bf Problem 2}. Construct a model in OpenMC for an infinite,
homogeneous system composed of a 1000-to-1 atom mixture of \isoA{H}{1} and
\isoA{U}{238}, with a point source of $1~\mathrm{MeV}$ neutrons. Compute the
neutron flux spectrum over $10^{-3}<E<10^5~\mathrm{eV}$. Then produce a plot
containing:
\begin{enumerate}[label=(\alph*)]
\item The computed flux spectrum.
\item The Maxwell-Boltzmann flux spectrum for the selected temperature, up to
      $1~\mathrm{eV}$.
\item The $1/E$ spectrum above $1~\mathrm{eV}$.
\end{enumerate}
Here the thermal comparison is a \emph{flux} spectrum,
$\phi_{\mathrm{MB}}(E)\propto E\exp[-E/(kT)]$, rather than the neutron-density
distribution used in Problem 1. The comparison curves may be scaled to compare
their shapes with the computed spectrum; no absolute source strength is
specified.

\begin{adaptedsolution}
One representative OpenMC model is
\begin{minted}{python}
import openmc
import numpy as np
import matplotlib.pyplot as plt

model = openmc.Model()
mat = openmc.Material(name="u238_h1")
mat.add_nuclide("H1", 1000.0)
mat.add_nuclide("U238", 1.0)
mat.set_density("g/cm3", 1.0)
mat.temperature = 294
model.materials = openmc.Materials([mat])

sphere = openmc.Sphere(r=1.0, boundary_type="reflective")
cell = openmc.Cell(region=-sphere, fill=mat)
model.geometry = openmc.Geometry([cell])

settings = openmc.Settings()
settings.run_mode = "fixed source"
settings.source = openmc.IndependentSource(
    space=openmc.stats.Point((0, 0, 0)),
    energy=openmc.stats.Discrete([1e6], [1.0]),
)
settings.seed = 1234
settings.batches = 100
settings.particles = 10000
settings.cutoff = {"energy_neutron": 1e-3}
model.settings = settings

energies = np.logspace(-3, 5, 1000)
tally = openmc.Tally()
tally.filters = [openmc.EnergyFilter(energies)]
tally.scores = ["flux"]
model.tallies = [tally]
\end{minted}
After running the model, plot the tally per unit energy and compare it with
scaled shapes:
\begin{minted}{python}
E = energies[:-1]
plt.step(E, flux_mean/np.diff(energies), where="post",
         color="k", lw=0.9, label="simulation")
plt.plot(E[E > 1], const_1/E[E > 1], "r--", label="1/E")
plt.plot(E[E < 1], const_2*E[E < 1]*np.exp(-E[E < 1]/0.0253),
         "b:", label="Maxwell-Boltzmann flux")
plt.xscale("log")
plt.yscale("log")
plt.xlabel("E [eV]")
plt.ylabel("flux per unit energy")
plt.legend()
\end{minted}
The simulated spectrum should approach a Maxwell-Boltzmann-like thermal shape
below about $1~\mathrm{eV}$ and a slowing-down shape close to $1/E$ above the
thermal range, with resonance effects where \isoA{U}{238} data are important.
\end{adaptedsolution}

\clearpage
